\documentclass[UTF8,a4paper]{ctexart}
% Chinese article in the style of the common Overleaf ctex templates:
% ctex picks XeLaTeX and the Fandol fonts; English and math stay in Latin fonts.
\usepackage{amsmath,amssymb}
\usepackage{booktabs}
\usepackage{graphicx}
\usepackage{xcolor}
\usepackage[colorlinks=true,linkcolor=blue]{hyperref}
\usepackage{geometry}
\geometry{margin=2.5cm}

\title{中文排版测试：混合中英文与数学公式}
\author{张三 \and 李四}
\date{\today}

\begin{document}
\maketitle
\begin{abstract}
本文测试中文排版的常见情况，包括标题、目录、中英文混排（Mixed Chinese and English）、
数学公式、表格以及\textbf{粗体}与\emph{强调}。
\end{abstract}

\tableofcontents

\section{引言}
\LaTeX{} 是一种高质量的排版系统。在中文环境下，我们通常使用 \texttt{ctex} 宏集。
标点符号如“引号”、《书名号》、（括号）以及省略号……都应正确处理。
繁體中文也應該可以正常顯示，例如：臺灣、香港。

\section{数学公式}
欧拉公式 $e^{i\pi} + 1 = 0$ 被称为最美的公式。高斯积分如下：
\begin{equation}
  \int_{-\infty}^{\infty} e^{-x^2}\,\mathrm{d}x = \sqrt{\pi}. \label{eq:gauss}
\end{equation}
由式~\eqref{eq:gauss} 可知，标准正态分布的密度函数积分为 $1$。

\section{表格}
\begin{table}[h]
  \centering
  \caption{各城市人口（单位：万）}
  \begin{tabular}{lrr}
    \toprule
    城市 & 2010 年 & 2020 年 \\
    \midrule
    北京 & 1961 & 2189 \\
    上海 & 2302 & 2487 \\
    广州 & 1270 & 1868 \\
    \bottomrule
  \end{tabular}
\end{table}

\section{列表}
\begin{enumerate}
  \item 第一项：{\heiti 黑体}与{\kaishu 楷体}；
  \item 第二项：{\fangsong 仿宋}；
  \item 第三项：\textcolor{red}{红色文字}。
\end{enumerate}

\end{document}
